\chapter{Introduction}
\label{ch:intro}

Approximate nearest neighbor search (ANNS) is a key operation in modern vector databases, retrieval systems, recommendation engines, and semantic search services.
At hundred-million scale and beyond, exact nearest neighbor search is usually too slow, so practical systems rely on approximate indexes such as locality-sensitive hashing\cite{datar2004lsh}, graph-based indexes\cite{malkov2020hnsw,fu2019nsg}, and quantization-based methods\cite{jegou2011pq}.
Among these approaches, graph-based ANNS provides a strong balance between recall and throughput\cite{subramanya2019diskann,chen2021spann}.

DiskANN is a representative disk-based graph ANNS system and has been adopted in practical vector databases and cloud services\cite{wang2021milvus,upreti2025cosmos}.
It combines the Vamana graph, product quantization (PQ), and beam search to support large indexes on a single machine while keeping memory usage manageable.
However, its search path still performs many random SSD reads.
Under concurrent serving, where many users issue retrieval requests at the same time, these reads accumulate in the shared SSD queue and become a major source of throughput loss and tail latency\cite{subramanya2019diskann}.

Prior work has improved disk-based ANNS from different directions.
For example, Starling improves page layout through block shuffling and an in-memory navigation graph, increasing the amount of useful information returned by each page read\cite{wang2024starling}.
This approach is effective under static-index settings, but its benefit depends on the offline layout and navigation-graph quality.
When memory budgets shrink, query distributions shift, or updates arrive over time, a fixed layout alone may not fully address the serving bottleneck.

This thesis proposes PaceANN, an acceleration layer for DiskANN that preserves the existing index layout.
PaceANN focuses on reducing unnecessary SSD I/O during search through two mechanisms.
First, the \textbf{Popularity-Adaptive Cache Manager} (PAC) dynamically caches frequently visited nodes using a sharded CLOCK design.
Each cache entry stores both the neighbor list and the full-precision vector, reducing I/O during traversal and re-ranking.
Second, \textbf{Convergence-Aware Early Termination} (CA-ET) adds a per-query stopping rule to beam search.
It uses already observed exact-distance candidates to decide whether the current query has converged, allowing easy queries to stop early while difficult queries continue searching\cite{li2020adaptiveet,chatzakis2025darth,teofili2025patience}.

The main contributions of this thesis are as follows:
\begin{itemize}
  \item We design PAC, a sharded CLOCK-based cache that adapts to serving-time access patterns and co-stores neighbor lists with full-precision vectors to reduce both traversal and re-ranking I/O.
  \item We introduce CA-ET, a per-query early-termination rule that detects convergence without issuing additional SSD reads.
  \item We show that PAC and CA-ET are complementary: caching stabilizes early search and supplies exact vectors, while early termination removes post-convergence expansions that caching alone cannot eliminate.
  \item We evaluate PaceANN under concurrent serving and equal memory budgets. On SPACEV100M, SIFT100M, and DEEP100M, PaceANN consistently outperforms DiskANN and Starling at the 100M scale, with QPS improvements of up to 59\%. Compared with PipeANN, the result is dataset-dependent: PaceANN leads on SPACEV100M, while PipeANN is stronger in the high-recall region on SIFT100M and DEEP100M.
\end{itemize}
